% Section 3.1: 01-metric-and-musical-isomorphisms | Part: 1 | Depends on: sec:vectors, sec:tensors | Status: done | Original lines: 787-833
%!TEX root = ../../../main.tex
% REWRITTEN: notation E -> V, m -> n; the sign conventions for the signature are stated once, in the pitfall of
% Section sec:quadform (the original sentence here pointed to sec:clifford, the pitfall to sec:ga-calculus).
\begin{formal}[{metric}]
A \emph{metric} on $V$ is a tensor $g\in\TT{0}{2}(V)$ (a bilinear map $g:V\times V\to\R$) such that
\begin{enumerate}[label=(\roman*),nosep]
\item $g(v,w)=g(w,v)$ (symmetric);
\item $g(v,w)=0\ \ \forall w\in V\ \Rightarrow\ v=0$ (non-degenerate).
\end{enumerate}
In a basis, $g=g_{ij}\,e^{i}\otimes e^{j}$ with $(g_{ij})$ a symmetric \emph{invertible} matrix. Being symmetric, it can be diagonalized over $\R$; if it has $p$ positive and $q$ negative eigenvalues, the \emph{signature} is $\sig(g)=(p,q)$ (it does not depend on the basis, by Sylvester's law of inertia). 
\begin{itemize}[nosep]
\item $\sig=(n,0)$ (positive definite): \textbf{Riemannian}: $g$ is an inner product.
\item $\sig=(n-1,1)$: \textbf{Lorentzian} (one negative direction, ``mostly plus''), used in GR with $n=4$, $\sig(g)=(3,1)$, \eg Minkowski $\eta=\mathrm{diag}(-1,1,1,1)$. The opposite (``mostly minus'') convention, $\eta=\mathrm{diag}(1,-1,-1,-1)$, and where each one is used, are listed in the pitfall of Section~\ref{sec:quadform}.
\end{itemize}
The \emph{norm squared} is $\|v\|^2=g(v,v)$, which can be positive, zero or negative in the Lorentzian case.
\end{formal}

\begin{intuition}[{the metric is the ``dot product'': it defines lengths, angles and (in GR) causality}]
$g(v,w)$ tells you the length of $v$, the angle between $v$ and $w$, and the projection of $v$ on $w$. In ordinary space all vectors have $g(v,v)\geq0$. In relativity the metric is \emph{indefinite}: a vector $v$ with $g(v,v)<0$ is \textbf{timelike} (a possible velocity of a massive particle), $g(v,v)=0$ with $v\ne0$ is \textbf{null} (light), $g(v,v)>0$ is \textbf{spacelike}. The null vectors form the \emph{light cone} and non-degeneracy does not mean ``$g(v,v)\neq0$ for all $v\neq0$'', it only means that $v\neq0$ always has \emph{some} partner $w$ with $g(v,w)\neq0$.
\end{intuition}

\begin{formal}[{musical isomorphisms (raising and lowering indices)}]
A metric gives an isomorphism between $V$ and $V^*$:
\[
g^\flat:V\to V^*,\quad v\mapsto g(v,\cdot),\qquad\qquad g^\sharp:=(g^\flat)^{-1}:V^*\to V.
\]
In components: $(g^\flat v)_k=v^jg_{jk}=:v_k$ (``lowering the index''), and, with $(g^{kl})$ the inverse matrix of $(g_{ij})$ (so $g^{ij}g_{jk}=\delta^i{}_k$), $(g^\sharp f)^l=f_kg^{kl}=:f^l$ (``raising the index''). The \emph{inverse metric} $g^{-1}=g^{ij}e_i\otimes e_j\in\TT{2}{0}$ is the induced metric on $V^*$. More generally the metric lets us raise or lower any index of any tensor, \eg $T^i{}_j\mapsto T_{ij}=g_{ik}T^k{}_j$; contraction with $g^{ij}$ is the way to take traces of $(0,2)$ tensors.
\end{formal}

\begin{intuition}[{the metric identifies arrows with rulers}]
Given an arrow $v$, the covector $g(v,\cdot)$ is ``the ruler that measures how much of any vector $w$ lies along $v$ (scaled by $|v|$)''. This is the \emph{only} reason why in Euclidean space we can pretend that arrows and covectors are the same thing: with $g=\delta_{ij}$ lowering and raising indices does not change the numbers. With another metric (polar coordinates, curved space, Minkowski) the numbers do change, and the difference between upper and lower indices matters.
\end{intuition}

\begin{example}[{Minkowski space}]
$\eta=\mathrm{diag}(-1,1,1,1)$, $\eta^{-1}=\mathrm{diag}(-1,1,1,1)$ too. For $v=(v^0,v^1,v^2,v^3)$ the lowered components are $v_\mu=(-v^0,v^1,v^2,v^3)$, so
\[
\eta(v,v)=-(v^0)^2+(v^1)^2+(v^2)^2+(v^3)^2.
\]
The vector $(1,1,0,0)$ is null (a light ray with $c=1$), $(2,1,0,0)$ is timelike, $(1,2,0,0)$ is spacelike. For a worldline $x^\mu(\tau)$, \emph{proper time} is $\dd\tau^2=-\eta(\dot x,\dot x)\,\dd\lambda^2$. A Lorentz transformation is precisely a basis change $A$ with $A^t\eta A=\eta$ (compare with $M'=A^tMA$ in Section~\ref{sec:tensors}).
\end{example}

% REORDERED (F2): the example "raising an index in polar coordinates" (original L826-831) moved to Section sec:pushpull.
